% page 336 of the facsimile (image p-335 of the scan)
% block 160.0 x 233.3 mm ; left margin 8.0 mm
\pgc{-12.700mm}{0.000mm}{
\l{\cel{5}328}
\l{}
\l{}
\l{\sou{lamo}.\cel{2}I. Peco de metalo, de ligno, de ivoro, e c., plata,}
\l{\cel{3}dina e streta. - II. Parto fera, stala, di instrumento,}
\l{\cel{3}di utensilo, apta tranchar, taliar, e c. - EFIS.}
\l{}
\l{\sou{lamao.}\cel{2}(\sou{zool.)}\cel{2}Quadripedo ruminera ek Peruvia, uzata kom}
\l{\cel{3}charjo-bestio, e di qua la pili uzesas por facar texuri.}
\l{\cel{3}- II. Sacerdoto di Buddha, en Tibet. - DEFIRS.}
\l{}
\l{\sou{lamantino}. (\sou{zool.})\cel{2}Cetaceo herbivora. - L.\cel{1}\sou{manatus}. DEFI.}
\l{}
\l{\sou{lamarckismo}. (\sou{biol.)}\cel{3}Teorio (transformista) da Lamarck, qua}
\l{\cel{3}asertas ke exkluzive la cirkonstanci extera efikas a la}
\l{\cel{3}materio vivoza. - DEFIRS.}
\l{}
\l{\sou{lamelo}. I. (\sou{tekniko)}\cel{2}Mikra lamo de vitro sur qua onu pozas}
\l{\cel{3}la objekto quan onu exploras per mikroskopo. - II.(\sou{bot.})}
\l{\cel{3}Apendico di la korolo ye qua konsistas la krono. - III.}
\l{\cel{3}(\sou{historio naturala}).Mikra lamo ek bloko kayeroza, folioza.}
\l{\cel{3}- DEFIS.}
\l{}
\l{\sou{lamentar}.\cel{2}(\sou{netrans., pri, pro}).\cel{2}Lasar explozar, lor la emoco}
\l{\cel{3}de sufro, klameti iterata. - DEFIS.}
\l{}
\l{\sou{laminar}.\cel{2}(\sou{trans.)}\cel{2}(\sou{tekn.)}\cel{2}Transformar (bloko metala) a lami,}
\l{\cel{3}per igar lu pasar inter du rotaco-cilindri, qui turnas}
\l{\cel{3}inversa-sinse. -II. Igar cilindro metala tre pezoza rular}
\l{\cel{3}sur (voyo) til kande ica divenabos glata de la ensuligeso}
\l{\cel{3}ed aplasteso di la stoni o stoneti qui saliis sur lu. -EFIS}
\l{}
\l{\sou{lampo}. - Utensilo konsistanta (origine) ye tanko qua kontenas}
\l{\cel{3}liquido kombustebla, en qua esas imersita mecho quan onu}
\l{\cel{3}acendas por produktar lumo. - Utensilo sama-skopo, en qua}
\l{\cel{3}lumo produktesas dafilo (metala o karba) inkandecoza per}
\l{\cel{3}elektro. - DEFIRS.}
\l{}
\l{\sou{lampaso}. Silko-stofo, quan onu obtenis olim de Chinia,kun}
\l{\cel{3}granda desegnuri reliefa, maxim-multa-kaze sur fondo altra-}
\l{\cel{3}kolora. - DEFIS.}
\l{}
\l{\sou{lampiono}. Kupeto qua kontenas materio kombustebla, kun mecho,}
\l{\cel{3}uzata por la ilumini. FI.}
\l{}
\l{\sou{lampiro}. (\sou{zool.}) Insekto di qua la femino esas sen-ala e fos-}
\l{\cel{3}forecas en nia klimato, -- aloza e fosforecoza pri la du}
\l{\cel{3}sexui, en la landi tropikala. -L.\cel{1}\sou{lampyris}.\cel{2}EF.}
\l{}
\l{\sou{lampredo.}\cel{2}(\sou{zool.)}\cel{2}Fisho ek la familio \textquotedbl{}ciklostomi\textquotedbl{}, di qua}
\l{\cel{3}la brankii esas sep-trua. - L.\cel{1}\sou{petromyzon}.. DEFIS.}
\l{}
\l{\sou{lano}.\cel{2}Pili dolca e tufatra, qui kreskas sur la pelo di la}
\l{\cel{3}mutoni, e c. - FIS.}
}
